\chapter{Einleitung}

Die letzten Jahre haben gezeigt, dass künstliche Intelligenz (KI) und insbesondere Chatbots
zunehmend Einzug in den Coaching-Sektor hatten und als digitale Selbstcoaching-Tools
betrachtet werden \citep{mai2023chatbots}. KI stellt eine Schnittmengen zur Datenwissenschaft dar, wobei durch den Einsatz von Techniken und Systemen intelligentes Verhaltens erzeugt wird\citep{balzert2009software}. 
Ein KI-gestützter Chatbot ist eine Softwareanwendung, die auf künstlicher Intelligenz basiert und mittels natürlicher Sprache mit Menschen kommuniziert. 
Sie werden als Large Language Models (LLMs) bezeichnet. Dabei handelt es sich um KI-Modelle, die auf der Verarbeitung und Erzeugung von natürlicher Sprache
basieren. Sie nutzen riesige Mengen an Textdaten, um statistische Muster zu lernen und
darauf basierend Outputs zu generieren \citep{Douglas2023Large}.
Als Designmerkmale werden gestaltungsbezogene Eigenschaften eines LLM Systems verstanden. Um abzuleiten, welche Ausprägung Designmerkmale haben müssen,
wird das Verständnis vom Anthropomorphismus betrachtet. Dies ist ein wechselseitiges Phänomen, bei dem Designer menschenähnliche Merkmale in Artefakte einbetten und die Betrachter ihre kognitive Reaktion auf diese Merkmale auf die Artefakte projizieren \citep{Xiao2025Humanizing}. Interaktionsmerkmale von Sprachmodellen sind die Eigenschaften der Mensch Modell Interaktion, die sich in der Gestalt der Prompts ausdrücken – also wie Aufgaben als Anweisungen formuliert, strukturiert und iterativ verfeinert werden, um das gewünschtes Verhalten des Modells zu erreichen. Ein Prompt ist eine natürliche Sprachinstruktion, mit der der Mensch das Modell steuert \citep{Gao2023Prompt}. 
Der Begriff Coaching bezeichnet eine zeitlich begrenzte, interaktive Begleitung, bei der Coaches Klienten dabei unterstützen, persönliche oder berufliche Ziele eigenverantwortlich zu definieren und zu erreichen. Der sprachliche Ursprung liegt im ungarischen Wort „kocsi“ (Kutsche) als Metapher für ein Beförderungsmittel ans gewünschte Ziel \citep{sonesh2015power, nyfoudi2021talent}.
Coaching hat sich in verschiedenen Praxisfeldern als personenorientiertes Beratungsformat etabliert, das die Selbststeuerung und
persönliche Entwicklung von Einzelpersonen, Teams und Organisationen unterstützt. Neben
klassischen Präsenzformaten gewinnen digitale Coaching-Formate, wie E-Mail-, Chat- oder
VR-Coaching, zunehmend an Bedeutung \citep{strong2020chatbots}. 
Das wesentliche Merkmal von VR ist die Inklusion, dass heißt, das Umgeben-Sein von einer Umgebung. VR versetzt den Teilnehmenden mitten in Informationen hinein \citep{bricken1990learning}. 
Vorwiegend KI gestützte Chatbots werden eingesetzt, um Soft Skills zu trainieren oder Reflexionsprozesse
zu unterstützen \citep{mai2020ai}. Softskills sind überfachlichen verhaltens- und beziehungsbezogenen sowie emotionale und kognitive Kompetenzen \citep{dell2017educational, werning2023virtual}.

\section{Problemstellung}

Im Jahr 2021 waren über 46,7 Millionen amerikanische Erwachsene von psychischen Erkrankungen betroffen, wobei nur 41,3\% im Laufe des Jahres eine psychotherapeutische
Behandlung erhielten. Mehr als 60\% derjenigen, die im Vorjahr behandelt wurden,
gaben an, virtuelle Angebote in Anspruch genommen zu haben. Als hinderliche Gründe psychotherapeutischer Behandlungen werden Kosten, Zeitaufwand, Unwissenheit über Anlaufstellen und die Annahme, die Probleme
selbst bewältigen zu können herausgestellt. Es besteht ein Zusammenhang zwischen psychischen Erkrankungen, der Inanspruchnahme von Behandlungen und einem
wahrgenommenen ungedeckten Bedarf an Versorgung \citep{Forman-Hoffman2025Prevalence}.
Trotz der wachsenden Verbreitung von Coaching-Chatbots bleibt unklar, in
welchem Umfang diese menschliche Therapie oder Coaching ersetzen können, insbesondere
in sensiblen Kontexten wie der Traumaarbeit. Im Folgenden wird der Begriff Trauma nach
dem Verständnis von \cite{Levine2010}  verwendet, da im Rahmen dieser Arbeit nicht die
klinische Diagnostik, sondern eine traumasensible Ausrichtung im Coaching-Kontext im
Vordergrund steht. Demnach wird Trauma nicht als Störung, sondern eine Verletzung,
verursacht durch Angst, Hilflosigkeit und Verlust verstanden. Levine beschreibt die damit
verbundene angeborene Fähigkeit zur Heilung durch Selbstregulation von Erregungszuständen und intensiven Emotionen. Traumasensibles Coaching ist demnach darauf ausgerichtet,
selbstregulatorische Prozesse zu unterstützen. Daraus ergeben sich zentrale Erfolgsfaktoren menschlichen Coachings wie Empathie, Wertschätzung und Vertrauen. Diese sind bislang nur begrenzt auf
die Mensch-Maschine-Interaktion übertragbar \citep{mai2020ai}.  Es bestehen ethische Bedenken und Risiken hinsichtlich fehlender emotionaler Sicherheit, unzureichender und unklarer Verantwortlichkeiten \citep{passmore2023can}.

\section{Zielsetzung}

Ein großes Thema bei der Entwicklung erfolgreicher Sprachmodelle ist, sich mit der Personalisierung dieser zu beschäftigen, um Interaktionen, Inhalte und Empfehlungen an die
individuellen Nutzerpräferenzen anzupassen \citep{Tan2024}. Bezüglich der Anwendung von LLMs im therapeutischen oder Coaching Kontext kann dies jedoch bedeuten, das individuelle, nicht entwicklungsfördernde Muster von Klienten vom System adaptiert werden \citep{Xiao2025Humanizing}.
Ziel der Arbeit ist es, ein theoretisch fundiertes Modell für ein KI-gestütztes
Coaching-System zu entwickeln, das auf einer klar definierten Entscheidungslogik beruht. Im
Zentrum steht die Übertragung psychologischer und neurobiologischer Wirkmechanismen
auf künstliche intelligente Systeme. Dabei soll aufgezeigt werden, wie dialogische KI gestaltet
werden kann, dass sie nicht lediglich bestehende Denk- und Verhaltensmuster reproduziert,
sondern zur aktiven Infragestellung und potenziellen Transformation dieser Muster beiträgt,
um auf diese Weise Entwicklungsprozesse zur Traumaheilung der Klienten voranzubringen.
Unter dialogischer KI werden Dialogsysteme verstanden, die über natürliche Sprache mit
Nutzenden interagieren und dabei nicht nur reaktiv auf Eingaben antworten, sondern durch
proaktive Impulse den Gesprächsverlauf aktiv mitgestalten. Dadurch können sie gezielt
Reflexionsprozesse anstoßen und Veränderung unterstützen \citep{Deng2025Proactive}. Die Untersuchung ist besonders relevant, da die zunehmende Digitalisierung der Coachingpraxis
neue Chancen, aber auch Herausforderungen für die professionelle Beratung und ethische
Standards birgt \citep{mai2023chatbots,passmore2023can}.



\section{Forschungsfrage}

\begin{quote}
Welche Design- und Interaktionsmerkmale sind erforderlich, damit ein KI-basierter Chatbot im Coaching-Kontext eine traumasensible Unterstützung zum Prozess der Traumaüberwindung für Personen mit leichten bis moderaten psychischen Belastungen gewährleisten kann?
\end{quote}

\subsection{Einordnung in bisherige Forschung}
Grundlage der Arbeit ist ein theoretisch-konzeptioneller Forschungsansatz, mit dem Ziel, ein strukturbasiertes Modell für den Einsatz von KI-gestützten Chatbots im Coachingkontext zu entwickeln. Ausgangspunkt ist somit einerseits fundiertes Wissen über die Funktionsweisen und Entwicklung künstlicher Intelligenz, insbesondere KI-basierter Sprachmodelle, sowie die Betrachtung psychologischer Theorien. 
Die Persönlichkeits-System-Interaktionen Theorie (PSI) von \cite{Kuhl2000The} erklärt motivationale und willentliche Phänomene wie Konzentration, den Umgang mit Misserfolgen, intrinsische Bindung an persönliche Ziele, Beharrlichkeit und die Umsetzung von Intentionen. Es werden Interaktionsmuster zwischen vier kognitiven Systemen – Denken und Intentionsgedächtnis, Fühlen und Erweiterungsgedächtnis, diskrepanzsensitive Objekterkennung und intuitive Verhaltenskontrolle – durch affektive Veränderungen moduliert. Dies äußert sich in einer veränderten Aktivierung der kognitiven Systeme und damit in beobachtbaren Veränderungen von Selbstregulation, kognitivem Fokus und Handlungsfähigkeit. Dabei unterstützen selbstregulatorische Fähigkeiten affektive Veränderungen beim Lernen und in der Therapie, während entwicklungsbedingte und erzieherische Risikofaktoren diese beeinträchtigen können. Menschliches Verhalten wird nicht von einzelnen Eigenschaften, sondern der Interaktion mehrerer psychischer Systeme beeinflusst. Die PSI-Theorie soll als neurobiologischer Ansatz die Grundlage des betrachteten Problems darstellen. 
Der Anwendungsrahmen der Arbeit bezieht sich auf Coachingprozesse im Kontext leichter bis mittelschwerer psychischer Belastungen bzw. stabiler, aber dysfunktionaler Verhaltensmuster. Leichte Belastungen sind durch milde, vorübergehende Symptome gekennzeichnet, während moderate Belastungen anhaltende und intensivere emotionale Reaktionen hervorrufen, ohne jedoch zu einer schweren Funktionsbeeinträchtigung zu führen. Leichte bis moderate Belastungen können physiologische und psychologische Veränderungen verursachen, die jedoch noch nicht das Ausmaß erreichen, das eine klinische Störung oder akute Gefährdung darstellt \citep{Ovsiannikova2024Peculiarities}.  
Ausgehend von diesem Ansatz untersucht die Arbeit:

    \begin{itemize}
    \item Wie Prompts gestaltet sein müssen, um diagnostisch und reflexionsfördernd zu wirken
  \item Welche Steuerungsmechanismen notwendig sind, um gezielt unterschiedliche psychische Systeme im Sinne der PSI-Theorie zu aktivieren
  \item Welche Risiken bei der Nutzung von KI-Sprachmodellen im traumasensiblen Coaching-Kontext existieren
    \end{itemize}

Zusammenfassend soll die Arbeit herausstellen, wie sich zentrale Mechanismen der PSI-Theorie in konkrete Design- und Interaktionsmerkmale eines KI-basierten Chatbots überführen lassen. 

\section{Aufbau der Arbeit}



